\documentclass[11pt]{article}
\usepackage[margin=27mm]{geometry}
\usepackage{amsmath,amssymb,amsthm,hyperref}
\newtheorem{lemma}{Lemma}
\newtheorem{theorem}[lemma]{Theorem}
\newcommand{\PSL}{\operatorname{PSL}}
\newcommand{\PGL}{\operatorname{PGL}}
\newcommand{\Aut}{\operatorname{Aut}}
\newcommand{\N}{\operatorname{N}}
\title{Full automorphisms of the dihedral subdivision family}
\author{Human Compression Project}\date{}
\begin{document}\maketitle
\begin{abstract}
Suppressing vertices of degree two reduces the
$(S_4,D_{16},D_8)$ incidence graph to a cubic graph. Tutte's bound on
cubic arc stabilizers then leaves at most an index-two extension of
the specified simple group. The action on the conjugacy class of a
vertex stabilizer excludes that extension. This proves the four
full-automorphism assertions for this family on p.~386 of the fixed
PSL$(2,q)$ source, conditional only on its separately verified subgroup
configurations and the classical inputs stated below.
\end{abstract}

\section{Suppression and the normalizer argument}
\begin{lemma}\label{suppress}
Let $X$ be a finite connected bipartite simple graph with degrees three
and two. Suppose a group $G$ preserves types, is edge transitive and
acts two-transitively on the three neighbors of each vertex of degree
three. If there are more than two vertices of degree three, suppressing
all vertices of degree two produces a connected simple cubic graph $Y$.
Moreover, $\Aut(X)=\Aut(Y)$ under their natural actions, and $G$ is
arc transitive on $Y$.
\end{lemma}
\begin{proof}
Replace each path $(u,e,v)$, with $e$ of degree two, by an edge $uv$.
There are no loops because $u,v$ are distinct neighbors in the original
simple graph. A priori parallel edges are possible. At any vertex $u$,
two of its three incident edges in the resulting multigraph may have
the same other endpoint. Having the same endpoint is a
$G_u$-invariant equivalence relation on those three edges. The
two-transitive action is primitive, so if such a pair exists all three
edges at $u$ have one common endpoint $v$. Cubicity at $v$ then makes
these two vertices a whole connected component. This contradicts
connectedness and the hypothesis on the number of vertices. Thus $Y$
is simple and cubic; connectedness is preserved by suppression.

Every automorphism of $X$ preserves the two degrees and hence induces
an automorphism of $Y$. Since there are no parallel edges, a vertex of
degree two is uniquely determined by its two neighbors, so the induced
map is injective. Conversely an automorphism of $Y$ extends uniquely
to its subdivision by its action on edges. This proves the asserted
group identification. Edge transitivity in $X$ gives transitivity on
incidences of vertices and edges in $Y$, which is arc transitivity.
\end{proof}

\begin{lemma}\label{normalizer}
Let a group $G$ act faithfully and transitively on a finite set $\Omega$,
and let its point stabilizer $M$ satisfy $\N_G(M)=M$.
If $G\triangleleft H\le\operatorname{Sym}(\Omega)$, conjugation
embeds $H$ into $\Aut(G)$. Its image preserves the $G$-conjugacy class
of $M$.
\end{lemma}
\begin{proof}
An element $c$ centralizing $G$ sends a chosen point $x$ to a point
$gx$ fixed by $M=G_x$. The two point stabilizers have the same order,
so $M=G_{gx}=gMg^{-1}$ and $g\in\N_G(M)=M$. Therefore $cx=x$.
Commutation with the transitive group now implies that $c$ fixes every
point, so $c=1$. The kernel of conjugation from $H$ to $\Aut(G)$ is
thus trivial. An element $h\in H$ takes $G_x$ to $G_{hx}$, and all
point stabilizers are $G$-conjugate. This proves the last assertion.
\end{proof}

\section{The specified family}
We use the following classical inputs as premises, not as new results.
For a finite connected simple cubic arc-transitive graph, Tutte's
theorem gives a vertex stabilizer of order at most $48$ in its full
automorphism group. Indeed that full group acts regularly on the
$s$-arcs for some $s\le5$, and there are $3\cdot2^{s-1}$ such arcs
from a fixed vertex. The two parts of this result are in
Tutte's 1947 and 1959 papers [2,3].

For an odd prime $p$ in the cases below,
$G=\PSL(2,p)$ is nonabelian simple,
$\Aut(G)=\PGL(2,p)$, and the two $G$-conjugacy classes of
subgroups isomorphic to $S_4$ are interchanged by every outer
automorphism. These are the classical subgroup and automorphism
facts used by the source; the last fact is explicitly recalled on
its p.~386. The subgroups $S_4$ are maximal, hence self-normalizing:
their normalizers could otherwise only be $G$, contradicting simplicity.

\begin{theorem}\label{aut}
Let $p\in\{17,31,79,97\}$. Suppose $M,L\le G=\PSL(2,p)$ give
one of the connected locally two-transitive coset graphs with
\[
 M\cong S_4,\qquad L\cong D_{16},\qquad M\cap L\cong D_8.
\]
Then the full automorphism group of its incidence graph is exactly the
given group $G$.
\end{theorem}
\begin{proof}
The two degrees are three and two. The half with stabilizer $M$ has
$[G:M]=p(p^2-1)/48>2$ vertices. Lemma~\ref{suppress} gives a connected
simple cubic graph $Y$ with the same full automorphism group $H$ and
with $G$ arc transitive. The action of $G$ on $Y$ is faithful, since
it is its action on $G/M$ and the core of the proper subgroup $M$ in
the simple group $G$ is trivial.

By Tutte's bound and orbit--stabilizer,
\[
 [H:G]=\frac{|H_x|}{|G_x|}\le\frac{48}{24}=2.
\]
Thus $G$ is normal in $H$, either trivially or as a subgroup of index
two. Lemma~\ref{normalizer} embeds $H$ into $\Aut(G)$, with image
preserving the conjugacy class of $M$. Every outer automorphism
interchanges the two $S_4$ classes, so this image contains no outer
automorphism. It follows that $H=G$. The equality transfers back to
the incidence graph by Lemma~\ref{suppress}.
\end{proof}

The proof applies to every configuration with these hypotheses; it does
not inspect a graph separately for each listed prime. More generally it
applies to any prime for which the same classical premises and subgroup
configuration hold. The count of configurations and their fusion under
outer automorphisms are separate assertions, treated by the subgroup
classification. The conclusion here supplies an upper bound on the full
graph group; merely observing that an outer automorphism interchanges
two configurations would not have supplied that bound.

The source's cubic $(S_4,S_4,D_8)$ graphs require a separate argument.
Their specified simple group has two vertex orbits, so Tutte's bound
cannot be applied to their full group without first establishing vertex
transitivity. We make no such inference in this note. No whole-source
completion, Lean formalization, new conjecture count or measured
compression factor is claimed by this component.

\begin{thebibliography}{9}
\bibitem{source} F.~Buekenhout, J.~De Saedeleer and D.~Leemans,
\emph{On the rank two geometries of the groups $\PSL(2,q)$: part II},
Ars Mathematica Contemporanea \textbf{6} (2013), 365--388, p.~386.
\bibitem{tutte47} W.~T.~Tutte, \emph{A family of cubical graphs},
Proc. Cambridge Philos. Soc. \textbf{43} (1947), 459--474,
\href{https://doi.org/10.1017/S0305004100023720}{doi:10.1017/S0305004100023720}.
\bibitem{tutte59} W.~T.~Tutte, \emph{On the symmetry of cubic graphs},
Canadian J. Math. \textbf{11} (1959), 621--624,
\href{https://doi.org/10.4153/CJM-1959-057-2}{doi:10.4153/CJM-1959-057-2}.
\end{thebibliography}
\end{document}
